% !TeX root = ../main.tex

\chapter{研究方法}

\section{系統硬體架構與光路設計}

\subsection{有限共軛光路設計 (Finite-Conjugate Optical Design)}
本研究之顯微鏡光路參考 OpenFlexure \citep{collins2020robotic} 之設計理念，採用**有限共軛 (Finite-Conjugate)** 光路系統，而非傳統研究級顯微鏡常用之無限遠校正 (Infinity-corrected) 系統。在此架構下，物鏡直接將樣本成像於感測器平面，省略了昂貴且佔空間的管鏡 (Tube lens) 與無限遠空間。此設計不僅大幅縮減了顯微鏡的物理體積，亦降低了光學組裝的複雜度，極其適合整合於樂高積木建構的可重構架構中。

感測端採用 Raspberry Pi Camera Module 2 (IMX219)，其原始鏡頭在移除後，配合特定設計的 3D 列印固定件，可與顯微物鏡形成穩定的光學鏈。雖然有限共軛系統在插入濾光片或分光鏡時會產生額外的球差，但本研究透過精確計算元件位置，將此影響控制在可接受範圍內。

\subsection{基於樂高與龍門架構之機械系統}
硬體支撐結構由樂高積木 (LEGO Technic系列) 構成之「光學底板」與客製化 3D 列印組件組成。為了解決傳統開源顯微鏡行程受限的問題，本系統開發了剛性龍門式 (Gantry-style) X-Y 掃描載台。
\begin{itemize}
    \item \textbf{動力系統}：採用三顆 NEMA 17 步進馬達分別驅動 X、Y、Z 三軸。
    \item \textbf{控制核心}：使用 Arduino Uno 搭配 CNC Shield V3 作為馬達驅動器，並與 Raspberry Pi 透過 Serial 通訊。
    \item \textbf{定位精度}：系統設定為 16 倍細分 (Microstepping)。根據 $Z$ 軸導螺桿導程為 2.0mm 且步進角為 1.8$^\circ$ (每圈 200 步) 計算，其理論解析度可達：
    \begin{equation}
        Resolution_Z = \frac{2000~\text{\textmu m}}{200 \times 16} = 0.625~\text{\textmu m/step}
    \end{equation}
    此精度足以滿足高倍率物鏡下的自動對焦需求。
\end{itemize}

\subsection{光源與電子控制}
光源系統包含明視野 (Bright-field) LED 與螢光激發光源。
\begin{itemize}
    \item \textbf{明視野}：由 Raspberry Pi GPIO 17 腳位控制，提供透射式照明。
    \item \textbf{螢光激發}：由 GPIO 27 腳位控制。透過 RPi.GPIO 函式庫實現與相機曝光時間的同步切換，避免不必要的樣本螢光淬滅 (Photobleaching)。
\end{itemize}

\section{軟體架構與運算成像管線}

\subsection{非同步多執行緒後端}
本系統之軟體基於 Python 撰寫，採用 PySide6 框架構建非同步多執行緒架構。為確保即時影像串流與硬體控制互不干擾，系統拆分為以下核心模組：
\begin{enumerate}
    \item \textbf{Camera Worker}：負責 Picamera2 的即時採集、緩衝管理與影像校正運算。
    \item \textbf{Motor Controller}：負責 G-code 指令的序列化發送與 Arduino 的回應監聽。
    \item \textbf{API Handler}：基於 HTTP 協定建立 Web 端控制介面，使使用者能遠端操作顯微鏡。
\end{enumerate}

\subsection{多指標動態自動對焦演算邏輯}
為應對明場與螢光影像特性的巨大差異，系統實作了「三重指標 (Triple-Metric)」動態搜尋演算法 (如圖 \ref{fig:af_logic}，概念圖)。
\begin{itemize}
    \item \textbf{第一階段 (粗調)}：使用全域變異數 (Global Variance)。其對離焦產生的模糊具備極佳的單調性，能引導馬達快速進入焦平面鄰域。
    \item \textbf{第二階段 (中調)}：切換至空間頻率 (Spatial Frequency)，利用二維梯度分佈縮小搜尋步進。
    \item \textbf{第三階段 (精調)}：在明場下使用 Brenner 梯度函數；在螢光下則啟動專屬的「形態學 + 前 0.1\% 梯度」算法。該算法先以形態學開運算移除背景霾霧 (Haze)，再提取最高能量的 0.1\% 像素計算梯度，確保在低光下仍能鎖定峰值。
\end{itemize}
若搜尋過程中指標連續下降，系統會判定進入「盲區」，自動觸發方向反轉或重設起始點，並具備「尋峰回溯」功能以修正馬達背隙造成的誤差。

\subsection{空間變異張量影像校正管線}
針對低成本感測器在顯微系統中的缺陷，本研究實作了高度優化的影像校正管線：
\begin{enumerate}
    \item \textbf{黑電平補償}：強制扣除感測器底噪 (Offset=64.0)，提升暗部細節對比。
    \item \textbf{平場校正 (Gain Maps)}：預先拍攝均勻白場，計算 R, G, B 三通道的二維空間增益圖，校正主光線角度 (CRA) 不匹配導致的暗角。
    \item \textbf{色彩解混 (Color Unmixing)}：針對螢光通道間的光譜串擾，引入預訓練之 $3 \times 3 \times H \times W$ 空間變異張量。利用愛因斯坦求和運算 (\texttt{np.einsum}) 對全圖進行像素級的解混校正：
    \begin{equation}
        I_{corrected}(x, y) = T(x, y) \cdot I_{raw}(x, y)
    \end{equation}
    其中 $T$ 為空間變異之色彩變換矩陣，能還原樣本真實的螢光光譜顏色。
\end{enumerate}

\section{高解析度圖案化照明傅立葉疊層重建 (piFP)}
硬體端在靠近樣本處放置一個具備隨機微結構的擴散片。藉由 X-Y 馬達帶動擴散片產生微小位移，使樣本受圖案化照明 (Pattern illumination) 照射。軟體端則實作迭代相位復原演算法，在每一輪迭代中同步更新擴散片傳輸函數與樣本之波前複振幅資訊。此技術透過頻域合成，將傳統受限於物鏡數值孔徑的高頻資訊還原，預期可將有效解析度提升 2 至 3 倍。